\section{无失真离散信源编码}
\warmth{0}\quad\whisper{这一节把前面的熵变成编码结论：若允许很长的分组，平均码率或平均码长可以逼近 \(H(X)\)。}

\begin{strand}
无失真离散信源编码的核心问题是：把长度为 \(L\) 的信源序列 \(X^L\)
映射成码字，再由码字恢复出 \(\hat X^L\)。如果分组足够长，高概率序列集中在典型集里，
于是只需要给典型序列分配码字，所需码字数量大约是 \(2^{L H(X)}\)。
\end{strand}

\block{编码模型}
\trigger{当题目出现 \(X^L\)、码字集合 \(\mathcal{U}^k\)、编码映射和译码映射时，先画这条链。}

\begin{definition}[分组信源编码]
设离散无记忆信源的长度 \(L\) 输出序列取值于 \(\mathcal{X}^L\)，码符号集合为
\(\mathcal{U}\)，每个码字长度为 \(k\)，则码字集合为 \(\mathcal{U}^k\)。编码映射和译码映射分别为
\[
f:\mathcal{X}^L\to \mathcal{U}^k,
\qquad
g:\mathcal{U}^k\to \mathcal{X}^L.
\]
整个过程可写为
\[
X^L \xrightarrow{\ f\ } U^k \xrightarrow{\ g\ } \hat X^L,
\]
其中 \(\hat X^L\) 是译码恢复出来的
\answerin[3cm]{信源序列}。
\end{definition}

\begin{remark}
若 \(k\) 为固定值，就是等长编码；若不同信源序列允许使用不同长度的码字，就是变长编码。
本节先看分组等长编码的基本极限。
\end{remark}

\block{什么叫无失真}
\begin{definition}[严格无失真与渐近无失真]
严格无失真要求对每个输入序列都有
\[
\hat X^L=X^L.
\]
更常用的渐近无失真用错误概率描述：
\[
P_e^{(L)}=P\{\hat X^L\neq X^L\}.
\]
若当 \(L\to\infty\) 时存在一族编码方案使
\[
P_e^{(L)}\to 0,
\]
则称可以实现渐近无失真编码。
\end{definition}

\begin{yourturn}
把“严格无失真”和“渐近无失真”的差别填出来：
\[
\text{严格无失真：}\hat X^L=\answerin[2cm]{X^L},
\qquad
\text{渐近无失真：}P_e^{(L)}\to\answerin[1.4cm]{0}.
\]
\workspace[2]
\end{yourturn}

\block{码率}
\begin{definition}[编码速率]
若码符号集合 \(\mathcal{U}\) 的大小为 \(M\)，每个信源分组长度为 \(L\)，码字长度为 \(k\)，
则编码速率定义为
\[
R=\frac{k}{L}\log_2 M
\]
bit/source symbol。它表示平均每个信源符号需要多少二进制位来表示。
\end{definition}

\begin{example}
若用二进制码字，\(M=2\)，则
\[
R=\frac{k}{L}.
\]
若每个 \(L=4\) 的信源分组用 \(k=6\) 个二进制码元表示，则码率为
\[
R=\frac{6}{4}=1.5
\]
bit/source symbol。
\end{example}

\begin{yourturn}
若码符号集合大小为 \(M=4\)，长度 \(L=5\) 的信源序列被编码成长度 \(k=6\) 的码字，则
\[
R=\frac{k}{L}\log_2 M
=\frac{6}{5}\log_2 4
=\answerin[1.6cm]{2.4}.
\]
\workspace[2]
\end{yourturn}

\block{熵给出的可达极限}
\begin{theorem}[离散无记忆信源无失真编码定理]
设离散无记忆信源 \(X\) 的熵为 \(H(X)\)。对任意 \(\epsilon>0\)，若编码速率满足
\[
R=\frac{k}{L}\log_2 M>H(X)+\epsilon,
\]
则当 \(L\) 足够大时，存在一族分组编码使错误概率
\[
P_e^{(L)}<\delta
\]
其中 \(\delta>0\) 可以任意小。反过来，若
\[
R<H(X),
\]
则不可能实现渐近无失真编码。
\end{theorem}

\begin{remark}
这个定理的严格证明要用典型序列和大数定律，后面讲无失真信源编码定理时再完整证明。
这里先保留结论和直觉：高概率序列集中在典型集里，典型集大小约为
\(2^{L H(X)}\)，所以只要码字数能覆盖典型集，错误概率就可以随 \(L\) 增大而变小。
\end{remark}

\block{典型序列的角色}
\begin{definition}[典型集]
对长度为 \(L\) 的信源序列，把概率接近 \(2^{-L H(X)}\) 的高概率序列放入典型集，
记作 \(A_\epsilon^{(L)}\)。当 \(L\) 足够大时，
\[
P\{X^L\in A_\epsilon^{(L)}\}>1-\delta,
\]
并且典型集大小大约满足
\[
\left|A_\epsilon^{(L)}\right|\leq 2^{L(H(X)+\epsilon)}.
\]
\end{definition}

\begin{strand}
典型集解释了为什么“不是所有 \(n^L\) 个序列都同等重要”。虽然长度为 \(L\) 的序列总数是
\(n^L\)，但实际概率质量几乎都集中在典型集里。编码器只需要重点照顾典型序列，
非典型序列造成的错误概率可以随着 \(L\) 增大而变小。
\end{strand}

\begin{yourturn}
用典型集语言补完整个可达性思路：
\[
P\{X^L\in A_\epsilon^{(L)}\}>1-\delta,
\qquad
\left|A_\epsilon^{(L)}\right|\leq \answerin[3.4cm]{2^{L(H(X)+\epsilon)}}.
\]
因此只要
\[
M^k\geq \left|A_\epsilon^{(L)}\right|,
\]
就有足够多码字覆盖典型序列，也就是
\[
R=\frac{k}{L}\log_2 M>\answerin[2.6cm]{H(X)+\epsilon}.
\]
\workspace[3]
\end{yourturn}

\block{变长编码：平均码长逼近熵}
\begin{strand}
等长编码把每个长度为 \(L\) 的信源序列都塞进同样长的码字；变长编码则让高概率序列用短码字，
低概率序列用长码字。它关心的不是固定码字长度 \(k\)，而是平均码长 \(\bar K\)。
\end{strand}

\begin{definition}[平均码长]
设长度为 \(L\) 的信源序列 \(x^L\) 被编码成 \(m\) 元码字，码长为 \(K(x^L)\)。
平均码长定义为
\[
\bar K=\sum_{x^L\in\Xcal^L}\pmf(x^L)K(x^L).
\]
因此
\[
\frac{\bar K}{L}
\]
表示每个信源符号平均需要的 \(m\) 元码符号数。
\end{definition}

\begin{theorem}[变长编码定理]
设 \(X\) 为离散无记忆信源，熵为 \(H(X)\)，码符号集合大小为 \(m\)。
对长度为 \(L\) 的信源扩展进行唯一可译变长编码时，任意唯一可译码都满足
\[
\frac{\bar K}{L}\geq \frac{H(X)}{\log_2 m}.
\]
另一方面，存在一族 \(m\) 元即时码，使
\[
\frac{\bar K}{L}<\frac{H(X)}{\log_2 m}+\frac{1}{L}.
\]
也就是说，通过增大分组长度 \(L\)，变长编码的每符号平均码长可以任意逼近熵下界。
\end{theorem}

\begin{proof}
先看下界。任意唯一可译码的码长必须满足 Kraft 不等式，因此它诱导的平均码长不可能低于
信源扩展熵对应的 \(m\) 元码符号数：
\[
\bar K\geq \frac{H(X^L)}{\log_2 m}.
\]
由于信源无记忆，
\[
H(X^L)=L H(X),
\]
所以
\[
\frac{\bar K}{L}\geq \frac{H(X)}{\log_2 m}.
\]

再看可达性。对每个序列取 Shannon 码长
\[
K(x^L)=\left\lceil -\log_m \pmf(x^L)\right\rceil,
\]
则
\[
K(x^L)<-\log_m \pmf(x^L)+1.
\]
对 \(x^L\) 求平均，得到
\[
\bar K<\frac{H(X^L)}{\log_2 m}+1
=\frac{L H(X)}{\log_2 m}+1.
\]
两边除以 \(L\)，即得
\[
\frac{\bar K}{L}<\frac{H(X)}{\log_2 m}+\frac{1}{L}.
\]
\end{proof}

\begin{yourturn}
复述变长编码定理的关键夹逼：
\[
\frac{H(X)}{\log_2 m}
\leq
\frac{\bar K}{L}
<
\answerin[3.2cm]{\frac{H(X)}{\log_2 m}+\frac{1}{L}}.
\]
当 \(m=2\) 时，\(\log_2 m=1\)，所以
\[
H(X)\leq \frac{\bar K}{L}<\answerin[2cm]{H(X)+\frac{1}{L}}.
\]
\workspace[2]
\end{yourturn}

\begin{definition}[编码效率]
把平均码长换算成 bit/source symbol：
\[
R=\frac{\bar K}{L}\log_2 m.
\]
编码效率定义为
\[
\eta=\frac{H(X)}{R}.
\]
当 \(R\) 越接近 \(H(X)\)，编码效率越接近 \(1\)。
\end{definition}

\begin{yourturn}
若采用二进制变长编码，且分组长度 \(L\) 足够大，使 \(1/L<\epsilon\)，则定理给出
\[
H(X)\leq \frac{\bar K}{L}<H(X)+\answerin[1.4cm]{\epsilon}.
\]
这说明二进制变长编码的平均码长可以逼近
\answerin[2cm]{\(H(X)\)}。
\workspace[2]
\end{yourturn}

\recall{为什么无失真离散信源编码的极限码率是 \(H(X)\)，而不是 \(\log n\)？}
